\section{Introduction}
Probabilistic forecasts are typically evaluated using a chosen scoring rule, such as log loss or the continuous ranked probability score (CRPS). However, different proper scoring rules emphasize different aspects of predictive quality and may lead to different conclusions about the same forecast. This raises a fundamental question: can we construct a predictor that performs well simultaneously across a broad class of scoring rules, rather than optimizing for a single one?

Omniprediction provides a principled framework to address this question. Introduced by \cite{gopalan2022omnipredictors}, omniprediction seeks a predictor that minimizes the worst-case loss over a class of loss functions and a class of comparator functions. Formally, given a class of losses $\mathcal{L}$ and a class of competitors $\mathcal{F}$, the goal is to construct a predictor $\hat{p}(t)$ that minimizes
\[
\sup_{S \in \mathcal{L},\, f \in \mathcal{F}} 
\frac{1}{T}\sum_{t=1}^T [S(\hat{p}(t), Y_t) - S(f(t), Y_t)].
\]

A key motivation for this objective arises from downstream decision-making. In many applications, different users interact with forecasts through their own utility (or loss) functions. If a user believes a predictive distribution $P$ and chooses an action that minimizes their expected loss, then this induces a proper scoring rule. Consequently, minimizing loss over all proper scoring rules leads to minimizing the decision loss for all users simultaneously, regardless of how their objectives differ.

More explicitly, let $u$ denote a user-specific utility (or loss) function and let $\mathcal{A}$ denote the action space. Given a predictive distribution $P$, a user selecting an action $a(P;u) \in \arg\min_{a \in \mathcal{A}} \mathbb{E}_{Y \sim P}[u(a,Y)]$ induces a proper scoring rule $S(P,Y) := u(a(P;u), Y)$. %, and omniprediction over all proper scoring rules can be equivalently formulated as minimizing the worst-case decision loss across all users:
% \[
% \sup_{\|u\|_2 \leq 1,\, f \in \mathcal{F}} 
% \mathbb{E}\!\left[u\big(a(p(X);u), Y\big)\right] 
% - 
% \mathbb{E}\!\left[u\big(a(f(X);u), Y\big)\right].
% \]

In many practical settings, however, working with full predictive distributions is unnecessary or computationally burdensome. As a natural stepping stone, we instead consider predictors that output a collection of quantile estimates. In this setting, users base their decisions on these quantile summaries rather than the full distribution. For a class of decision problems in which optimal actions depend only on certain quantiles, a user who trusts the reported quantile estimates selects an action that is optimal with respect to those quantiles. {\color{blue} example to be added}

% A canonical example arises in interval prediction. A decision-maker selects lower and upper bounds $(a_\ell, a_u)$ and incurs asymmetric penalties for under- and over-coverage. Consider a utility of the form
% \[
% u(a_\ell, a_u, y)
% =
% c_\ell^- (a_\ell - y)_+ + c_\ell^+ (y - a_\ell)_+
% +
% c_u^- (a_u - y)_+ + c_u^+ (y - a_u)_+
% + \lambda (a_u - a_\ell),
% \]
% where $(x)_+ = \max(x,0)$ and $\lambda \ge 0$ controls the width of the interval.


Analogously to the probabilistic case, this decision process induces a class of scoring functions that are \emph{consistent} for the corresponding quantile functionals. Therefore, restricting attention to quantile forecasts leads naturally to an omniprediction objective over all scoring functions that are consistent for the specified quantile levels. This provides a principled bridge between the general probabilistic forecasting framework and the quantile-based setting studied in this paper.

While the omniprediction framework is conceptually appealing, its algorithmic realization is challenging. Recent work has made significant progress in special cases. For binary outcomes, \cite{kleinberg2023ucal} worked on minimizing the regret with respect to a single hypothetical base rate forecaster over all proper scoring rules. \cite{isaac2025omni} proposes two-player game-based algorithms that achieve near-optimal regret guarantees while competing against all proper scoring rules. For more general settings, \cite{gopalan2022omnipredictors} studies omniprediction in a batch framework and characterizes the statistical complexity of the problem via approximation-theoretic quantities. More recently, \cite{lu2025omni_nonlinear} considers omniprediction for non-linear losses and continuous outcomes, focusing on sample-efficient algorithms and swap regret guarantees.

Despite these advances, there remains a gap between theory and practice in continuous-outcome settings. In particular, many applications in forecasting focus on quantiles rather than full distributions, and require producing multiple quantile estimates that are coherent (i.e., non-crossing). Existing omniprediction approaches do not directly address this structure, nor do they provide efficient algorithms that exploit the rich geometry of quantile-based losses.

\paragraph{Our contributions.}
In this paper, we develop an omniprediction framework tailored to quantile forecasting, covering both single and multiple quantile levels. Our main contributions are as follows:

\begin{itemize}
    \item \textbf{Omniprediction for quantile functionals.} We formulate omniprediction over the class of all scoring functions that are consistent for quantiles, leveraging classical decomposition results that express such losses as weighted combinations of elementary threshold-based scores.
    
    \item \textbf{Efficient online algorithms.} We propose computationally efficient, two-player game-based algorithms for both single and multiple quantile levels. By exploiting the structure of consistent scoring functions, we derive closed-form solutions to the underlying minimax problems, avoiding the need for generic optimization.
    
    \item \textbf{Theoretical guarantees.} We establish regret bounds for our algorithms that scale as $\tilde{\mathcal{O}}(\sqrt{\log(mB)/T})$ for single quantile prediction and $\tilde{\mathcal{O}}(\sqrt{\log(mNB)/T})$ for multiple quantiles, matching standard online learning rates up to logarithmic factors.
    
    \item \textbf{Coherent multi-quantile prediction.} Our multiple quantile algorithm produces non-crossing quantile estimates by construction, even when base forecasts exhibit quantile crossing. We show that naively running single-quantile omniprediction independently fails to achieve this property and can be suboptimal in terms of the underlying minimax objective.
    
    \item \textbf{Extensions to weighted quantile losses.} We further study omniprediction under weighted sums of pinball losses and propose two algorithms with different computational and statistical trade-offs, along with a characterization of when post-processing (such as sorting) is beneficial.
\end{itemize}

\paragraph{Related work.}
Our work builds on and extends several lines of research. The omniprediction framework was introduced by \cite{gopalan2022omnipredictors}, who studied batch learning guarantees and connections to convex function approximation. The binary setting was further developed by \cite{isaac2025omni}, who proposed online algorithms with strong regret guarantees. In contrast, our work focuses on continuous outcomes and exploits the structure of quantile-based losses to obtain efficient algorithms.

The recent work of \cite{lu2025omni_nonlinear} considers omniprediction for general non-linear losses and emphasizes sample efficiency and swap regret. Their approach is more general but does not directly leverage the structure of quantile functionals or address coherence constraints such as non-crossing. Our approach is complementary: by specializing to quantile losses, we obtain simpler algorithms with explicit solutions and additional structural guarantees.

Finally, our work is closely related to the literature on proper scoring rules and consistent loss functions \cite{gneiting2011point, ehm2016binarydecomp, fissler2016multiq}, which provides the key decomposition tools that underpin our algorithms.

\paragraph{Organization of the paper.}
The remainder of the paper is organized as follows. In Section~\ref{sec:preliminaries}, we review background on proper scoring rules, consistent scoring functions, and the omniprediction framework. Section~\ref{sec:single_q} introduces our algorithm for single quantile omniprediction and establishes its theoretical guarantees. Section~\ref{sec:multi_q} extends the approach to multiple quantile levels and presents the corresponding algorithm and analysis. Section~\ref{sec:multi_wql} studies omniprediction for weighted quantile losses and compares alternative algorithms. Additional technical details and proofs are provided in the appendix.